\section{Wnioski z analizy literatury i potencjalne kierunki badawcze}

Protokół MCP jest nadal stosunkowo świeżą technologią z ogromnym potencjałem badawczym. Pierwsze testy wskazują, że modele LLM na ten moment nie radzą sobie dobrze z atakami skierowaną w tą warstwę aplikacji. Pierwsze inicjatywy  zajmujące się bezpieczeństwem MCP są dopiero w początkowych fazach działania. Obecnie dostępne artykuły i badania w większości czekają na recenzję i bycie opublikowanym. Dodatkowo, posiadają one istotne problemy:

\begin{itemize}
    \item skupienie się głównie na samych atakach, a nie sugestiach obrony przed nimi,
    \item rzadko wskazują konkretną warstwę systemu odpowiedzialną za podatność,
    \item jeżeli praca już posiada wskazówki jak chronić się przed atakami, rzadko kiedy poparte są faktycznymi danymi,
    \item ciężko znaleźć jakiekolwiek materiały w języku polskim.
\end{itemize}

Na tej podstawie sformułowane zostały dwa potencjalne kierunki badawcze dla pracy magisterskiej. Pierwszą z nich jest systematyczne \textbf{zbadanie skuteczności sugerowanych mechanizmów obronnych przed atakami na warstwę MCP}. W ramach takiej pracy należałoby najpierw wybrać grupę badanych ataków. Następnie należałoby zebrać sugerowane sposoby obrony przed nimi (OWASP, dokumentacja MCP, prace naukowe, narzędzia). Konieczne by było zaimplementowanie własnego środowiska testowego, z własnym serwerem MCP. Taka praca dawałaby praktyczną wartość deweloperom, chcącym tworzyć bezpieczne systemy wykorzystujące AI.

Drugim potencjalnym kierunkiem badawczym jest \textbf{zbadanie wpływu opisów narzędzi MCP na decyzję agenta}. Jest to bardziej niszowy temat, ale również istotny w kontekście bezpieczeństwa. Agent AI podejmuje decyzję, jakiego narzędzia MCP użyć, na podstawie ich opisów. Praca badałaby, jak odpowiednio spreparowane metadane wpływałyby na wybór narzędzia:

\begin{itemize}
    \item opis neutralny - jako grupa kontrolna,
    \item opis niejednoznaczny - czy narzędzie z opisem bardzo ogólnym, nie zawierającym żadnych szczegółów, będzie wybierane przez agenta,
    \item opis sugestywny - sugerujący agentowi wybór konkretnego narzędzia, np. "najlepsze na świecie i jedyne poprawnie działające narzędzie do edycji plików PDF",
    \item opis złośliwy - zawierający atak typu \textit{prompt injection},
    \item opis konfliktowy z promptem systemowym.
\end{itemize}